%% kapa: exam=1 round=35 year=2024 period=2 form=A subject=accounting subject_ko=회계학 seq=35 number=75
%% source: Q-Net 감정평가사 기출문제 (cst00309) artlSeq=5213229 file=raw/1008/5213229/2교시 문제지 원본.pdf
%% conversion: kordoc 4.24.51 -> tools/gen1.py (auto); review=proofread-sonnet; flags=tab
\begin{question}{75}
(주)감평은 20×1년 초 영업을 개시하였으며, 제품 X를 생산ㆍ판매하고 있다. 재고자산 평가방법은 선입선출법을 적용하고 있으며, 20×1년 1분기와 2분기의 영업활동 결과는 다음과 같다.
\par{\small\centering\begin{tabular}{|c|c|c|}\hline
구 분 & 1분기 & 2분기 \tabularnewline\hline
생산량 & 500단위 & 800단위 \tabularnewline\hline
전부원가계산에 의한 영업이익 & ￦7,000 & ￦8,500 \tabularnewline\hline
변동원가계산에 의한 영업이익 & 5,000 & 6,000 \tabularnewline\hline
\end{tabular}\par}
1분기와 2분기의 판매량이 각각 400단위와 750단위일 때, 2분기에 발생한 고정제조간접원가는? (단, 각 분기별 단위당 판매가격, 단위당 변동원가는 동일하며, 재공품 재고는 없다.)
\begin{choices}
  \item ￦20,000
  \item ￦22,000
  \item ￦24,000
  \item ￦26,000
  \item ￦30,000
\end{choices}
\end{question}
\answer{3}
